% Options for packages loaded elsewhere
\PassOptionsToPackage{unicode}{hyperref}
\PassOptionsToPackage{hyphens}{url}
\documentclass[
  ignorenonframetext,
]{beamer}
\newif\ifbibliography
\usepackage{pgfpages}
\setbeamertemplate{caption}[numbered]
\setbeamertemplate{caption label separator}{: }
\setbeamercolor{caption name}{fg=normal text.fg}
\beamertemplatenavigationsymbolsempty
% remove section numbering
\setbeamertemplate{part page}{
  \centering
  \begin{beamercolorbox}[sep=16pt,center]{part title}
    \usebeamerfont{part title}\insertpart\par
  \end{beamercolorbox}
}
\setbeamertemplate{section page}{
  \centering
  \begin{beamercolorbox}[sep=12pt,center]{section title}
    \usebeamerfont{section title}\insertsection\par
  \end{beamercolorbox}
}
\setbeamertemplate{subsection page}{
  \centering
  \begin{beamercolorbox}[sep=8pt,center]{subsection title}
    \usebeamerfont{subsection title}\insertsubsection\par
  \end{beamercolorbox}
}
% Prevent slide breaks in the middle of a paragraph
\widowpenalties 1 10000
\raggedbottom
\AtBeginPart{
  \frame{\partpage}
}
\AtBeginSection{
  \ifbibliography
  \else
    \frame{\sectionpage}
  \fi
}
\AtBeginSubsection{
  \frame{\subsectionpage}
}
\usepackage{iftex}
\ifPDFTeX
  \usepackage[T1]{fontenc}
  \usepackage[utf8]{inputenc}
  \usepackage{textcomp} % provide euro and other symbols
\else % if luatex or xetex
  \usepackage{unicode-math} % this also loads fontspec
  \defaultfontfeatures{Scale=MatchLowercase}
  \defaultfontfeatures[\rmfamily]{Ligatures=TeX,Scale=1}
\fi
\usepackage{lmodern}
\ifPDFTeX\else
  % xetex/luatex font selection
\fi
% Use upquote if available, for straight quotes in verbatim environments
\IfFileExists{upquote.sty}{\usepackage{upquote}}{}
\IfFileExists{microtype.sty}{% use microtype if available
  \usepackage[]{microtype}
  \UseMicrotypeSet[protrusion]{basicmath} % disable protrusion for tt fonts
}{}
\makeatletter
\@ifundefined{KOMAClassName}{% if non-KOMA class
  \IfFileExists{parskip.sty}{%
    \usepackage{parskip}
  }{% else
    \setlength{\parindent}{0pt}
    \setlength{\parskip}{6pt plus 2pt minus 1pt}}
}{% if KOMA class
  \KOMAoptions{parskip=half}}
\makeatother
\setlength{\emergencystretch}{3em} % prevent overfull lines
\providecommand{\tightlist}{%
  \setlength{\itemsep}{0pt}\setlength{\parskip}{0pt}}
% Slide layout defaults for warning-clean scientific decks.
\usepackage{etoolbox}
\IfFileExists{xurl.sty}{\usepackage{xurl}}{}
\IfFileExists{seqsplit.sty}{\usepackage{seqsplit}}{\newcommand{\seqsplit}[1]{#1}}
\protected\def\breakseq#1{\seqsplit{#1}}
\protected\def\breaktt#1{\begingroup\ttfamily\seqsplit{#1}\endgroup}
\setlength{\emergencystretch}{6em}
\tolerance=5000
\hbadness=10000
\hfuzz=1pt
\setlength{\tabcolsep}{2pt}
\AtBeginEnvironment{longtable}{\tiny\renewcommand{\arraystretch}{0.86}\setlength{\tabcolsep}{1pt}}
\AtBeginEnvironment{tabular}{\tiny\renewcommand{\arraystretch}{0.86}\setlength{\tabcolsep}{1pt}}

% Natbib and cross-reference fallbacks — slides don't load natbib
% or cleveref, but combined-PDF manuscript prose may emit these
% commands. The fallback renders citations as a bracketed key list
% and cross-references as detokenized labels so slides don't fail on
% undefined control sequences or raw underscores. \providecommand is
% a no-op if packages load later.
\providecommand{\citep}[1]{[#1]}
\providecommand{\citet}[1]{#1}
\providecommand{\citealp}[1]{#1}
\providecommand{\citeauthor}[1]{#1}
\providecommand{\citeyear}[1]{#1}
\providecommand{\cref}[1]{\texttt{\detokenize{#1}}}
\providecommand{\Cref}[1]{\texttt{\detokenize{#1}}}

\newtheorem{proposition}{Proposition}
\newtheorem{hypothesis}{Hypothesis}
\usepackage{bookmark}
\IfFileExists{xurl.sty}{\usepackage{xurl}}{} % add URL line breaks if available
\urlstyle{same}
% fallback for those not using the hyperref driver hyperxmp:
\makeatletter
\@ifundefined{xmpquote}{\newcommand{\xmpquote}[1]{#1}}{}
\makeatother
\hypersetup{
  hidelinks,
  pdfcreator={LaTeX via pandoc}}

\author{\texorpdfstring{}{}}
\date{}

\begin{document}

\section{Abstract}\label{abstract}

\begin{frame}[fragile,allowframebreaks]{Abstract}
Research software repositories in monorepo configurations accumulate
three categories of shared resources that individual projects must
consume without re-implementing discovery logic: \textbf{data pools}
(bibliographies, contacts, datasets), \textbf{governance rules} (style
guides, coverage thresholds, citation schemas), and \textbf{executable
tools} (code executors, validators, skill invocations). Without a
canonical integration pattern, projects either duplicate discovery logic
or silently ignore resources that fail to load --- both outcomes degrade
reproducibility and collaborative cohesion \breakseq{[@Wilson2014best;
@Taschuk2017ten]}.

This paper presents \breaktt{template\_pools\_rules\_tools}, a
meta-project exemplar that demonstrates how a single project can
programmatically discover, validate, and exercise all three resource
categories with zero tight coupling to any specific resource instance.
The exemplar comprises eight Python modules --- three resource readers
(\texttt{fonds\_reader}, \texttt{rules\_applier},
\texttt{tools\_invoker}), an orchestrator (\texttt{integration}), a
semantic rule evaluator (\breaktt{strong\_rule\_evaluator}), a figure
generator (\texttt{figures}), a manuscript-token generator
(\breaktt{manuscript\_variables}), and shared type definitions
(\texttt{type\_defs}) --- plus six thin orchestration scripts and a
fully token-injected manuscript pipeline.

The architecture (@fig:architecture) separates \emph{resource ownership}
from \emph{resource consumption}. Resources live in top-level
\texttt{fonds/}, \texttt{rules/}, and \texttt{tools/} directories and
are never modified by consumers. Each resource exposes a typed manifest
(\texttt{fonds.yaml}, \texttt{rules.yaml}, \texttt{tools.yaml}) that the
corresponding reader module uses for discovery and validation. All
readers implement graceful fallbacks: they return \texttt{None} or empty
collections when a resource is absent, log a warning via the standard
library \texttt{logging} module, and allow the integration pipeline to
continue. This revision extends the original three-figure presentation
to eight content figures plus a cover illustration --- a fond taxonomy
(@fig:taxonomy), a rule hierarchy (@fig:rulehier), a tool invocation
contract (@fig:toolcontract), a three-level resilience diagram
(@fig:resilience), and a script pipeline flow (@fig:pipelineflow) --- so
that every structural claim in the prose has a corresponding visual.

In a representative pipeline run, the integration demo loaded 3 fonds,
validated 2 rule sets, discovered 3 tools, and processed 8 bibliography
entries --- all reported as structured JSON that populates manuscript
variable tokens at render time. Tests covering the eight \texttt{src/}
modules (across nine test files) achieve well above the required ≥90\%
combined line coverage and use real file paths rather than mocks,
ensuring that reported counts are genuine --- run
\breaktt{uv\ run\ pytest\ …\ -\/-cov-report=term} for the current test
count and coverage percentage rather than trusting a number printed
here.

The \breaktt{template\_pools\_rules\_tools} exemplar provides a reference
implementation that any project in the template repository can consult
when designing its own resource-consumption layer.

This reference implementation is deliberately reproducible end-to-end:
every count, figure, and page-layout choice in this document is
regenerated from source rather than hand-authored, per the
Reproducibility Checklist in the front matter.

\textbf{Keywords:} research software engineering, monorepo architecture,
reproducibility, fonds, governance rules, tool discovery, graceful
degradation
\end{frame}

\end{document}
